% Chapter 3: Methodology
% Presents the research design framework, ontology construction process, 
% system development approach, and evaluation methodology
% Aligned with CRISP-DM framework from research proposal
\label{chap:methodology}

\section{Research Design Overview}
\label{sec:research_design}

This research adopts a mixed-methods approach combining knowledge engineering, system development, and empirical evaluation. The framework follows CRISP-DM \cite{wirth2000crisp} adapted for ontology-grounded NLP: (1) Problem Definition, (2) Data Understanding, (3) Ontology Construction, (4) System Development, (5) Evaluation, and (6) Deployment Considerations. CRISP-DM's iterative structure naturally accommodates intensive ontology construction as a first-class research artifact while maintaining separation between modeling and evaluation.

Figure~\ref{fig:methodology_overview} illustrates the adapted CRISP-DM workflow for this project, highlighting the six phases, key activities, deliverables, and iterative feedback loops that enabled systematic refinement of both the cultural ontology and the OG-RAG system.

\begin{figure}[htbp]
\centering
\resizebox{0.95\textwidth}{!}{% Methodology Flowchart - Figure 3.1
\begin{tikzpicture}[
    node distance=1.2cm and 2cm,
    phase/.style={rectangle, draw, fill=blue!20, text width=3.5cm, minimum height=1.1cm, align=center, rounded corners, font=\small, thick},
    activity/.style={rectangle, draw, fill=green!15, text width=3cm, minimum height=0.8cm, align=center, font=\footnotesize},
    arrow/.style={->, >=stealth, thick},
    feedback/.style={->, >=stealth, dashed, red, thick},
]

% Column 1: CRISP-DM Phases
\node[phase] (p1) at (0,0) {\textbf{1. Problem\\Definition}};
\node[phase, below=of p1] (p2) {\textbf{2. Data\\Understanding}};
\node[phase, below=of p2] (p3) {\textbf{3. Ontology\\Construction}};
\node[phase, below=of p3] (p4) {\textbf{4. System\\Development}};
\node[phase, below=of p4] (p5) {\textbf{5. Evaluation}};
\node[phase, below=of p5] (p6) {\textbf{6. Deployment}};

% Column 2: Key Activities
\node[activity, right=of p1] (a1) {Expert consultations\\Success criteria};
\node[activity, right=of p2] (a2) {Corpus analysis\\Theme extraction};
\node[activity, right=of p3] (a3) {847 concepts\\Neo4j schema};
\node[activity, right=of p4] (a4) {OG-RAG system\\3 baselines};
\node[activity, right=of p5] (a5) {3 annotators\\9 metrics};
\node[activity, right=of p6] (a6) {Documentation\\Framework};

% Column 3: Deliverables
\node[activity, fill=orange!15, right=of a1] (d1) {Research gap\\Hypotheses};
\node[activity, fill=orange!15, right=of a2] (d2) {100 proverbs\\Themes};
\node[activity, fill=orange!15, right=of a3] (d3) {KG: 947 nodes\\1,247 edges};
\node[activity, fill=orange!15, right=of a4] (d4) {thiLLMo\\system};
\node[activity, fill=orange!15, right=of a5] (d5) {Results:\\$p<0.001$};
\node[activity, fill=orange!15, right=of a6] (d6) {Thesis\\Code};

% Main workflow arrows
\draw[arrow] (p1) -- (p2);
\draw[arrow] (p2) -- (p3);
\draw[arrow] (p3) -- (p4);
\draw[arrow] (p4) -- (p5);
\draw[arrow] (p5) -- (p6);

% Horizontal connections
\draw[arrow] (p1) -- (a1);
\draw[arrow] (p2) -- (a2);
\draw[arrow] (p3) -- (a3);
\draw[arrow] (p4) -- (a4);
\draw[arrow] (p5) -- (a5);
\draw[arrow] (p6) -- (a6);

\draw[arrow] (a1) -- (d1);
\draw[arrow] (a2) -- (d2);
\draw[arrow] (a3) -- (d3);
\draw[arrow] (a4) -- (d4);
\draw[arrow] (a5) -- (d5);
\draw[arrow] (a6) -- (d6);

% Feedback loops (iterative refinement)
\draw[feedback] (p5.west) -- ++(-0.8,0) |- (p3.west) node[pos=0.25, left, font=\tiny, text width=1.2cm, align=right] {Error\\analysis};
\draw[feedback] (p4.west) -- ++(-0.5,0) |- (p3.west);
\draw[feedback] (p3.west) -- ++(-0.5,0) |- (p2.west) node[pos=0.25, left, font=\tiny, text width=1.2cm, align=right] {Corpus\\expansion};

% Annotations
\node[above=0.2cm of p1, font=\footnotesize\bfseries, text=purple] {CRISP-DM Adapted};
\node[above=0.2cm of a1, font=\footnotesize\bfseries, text=green!50!black] {Activities};
\node[above=0.2cm of d1, font=\footnotesize\bfseries, text=orange!80!black] {Outputs};

\end{tikzpicture}}
\caption{CRISP-DM Methodology Adapted for thiLLMo: Six-phase research workflow showing progression from problem definition through deployment. Each phase (blue) connects to specific activities (green) producing concrete deliverables (orange). Red dashed arrows indicate iterative refinement cycles: evaluation results fed back to ontology construction, system development informed ontology structure, and corpus analysis triggered data expansion. This adaptation emphasizes ontology construction as a first-class research phase rather than a preprocessing step.}
\label{fig:methodology_overview}
\end{figure}

\section{Phase 1: Problem Definition}
\label{sec:phase_business}

Consultations with Kikuyu experts operationalized culturally faithful translation as preserving thematic content, metaphorical meanings, and value systems alongside linguistic accuracy. Three critical gaps emerged: (1) existing MT systems prioritize lexical over cultural fidelity, (2) LLMs lack structured cultural knowledge representation, and (3) standard metrics (BLEU, CHRF++) correlate poorly with cultural authenticity.

\subsection{Success Criteria Definition}
\label{subsec:success_criteria}

Success criteria were established across three dimensions:

\begin{description}
    \item[Technical Performance] The OG-RAG system must demonstrate statistically significant improvements over baseline approaches (Traditional RAG, Raw LLM) on composite quality metrics combining cultural authenticity and translation fidelity. Target: $p < 0.05$ across all pairwise comparisons, with effect sizes (Cohen's $d$) exceeding 0.3 (medium effect).
    
    \item[Ontology Quality] The constructed Kikuyu proverb ontology must achieve: (1) structural validity (no critical pitfalls detected by OOPS! \cite{poveda2014oops}), (2) expert validation (approval by native Kikuyu speakers with cultural competence), and (3) sufficient coverage (capturing primary themes in wealth-related proverbs from Ireri's collection \cite{ireri2019proverbs}).
    
    \item[Practical Utility] The system must be computationally feasible (average translation latency $<$ 10 seconds) and produce outputs that native speakers judge to be culturally appropriate and pedagogically useful.
\end{description}

These criteria operationalize the broader research objectives from Section~\ref{sec:objectives}, making them empirically testable.

\section{Phase 2: Data Understanding}
\label{sec:phase_data_understanding}

The primary corpus comprised Ireri's 100 Kikuyu wealth proverbs \cite{ireri2019proverbs}, collected through elder interviews (2018-2019) across three counties. Supplementary sources included Gikandi \cite{gikandi2005thousand} and Kenyatta \cite{kenyatta1938facing}. Corpus analysis revealed: mean length 7.3 words (SD=2.8), 64\% requiring cultural interpretation, 12\% dialectical variation. Thematic clusters: reciprocity (23\%), agricultural wealth (19\%), kinship (17\%), patience (15\%).

\section{Phase 3: Ontology Construction}
\label{sec:phase_ontology}

The ontology was developed through iterative refinement following Noy \& McGuinness \cite{noy2001ontology}: (1) scope bounded to wealth-domain proverbs, (2) reused CIDOC CRM \cite{doerr2003cidoc} and SKOS templates, (3) extracted 847 terms via expert elicitation and corpus analysis, (4) organized into 5 top-level classes (Proverb, CulturalTheme, UsageContext, MoralLesson, CulturalEntity) with hierarchical sub-classes, (5) defined relationships (\texttt{expresses}, \texttt{teachesLesson}, \texttt{usedIn}) and properties, (6) instantiated 100 proverbs with manual annotation, (7) validated via OOPS! \cite{poveda2014oops}, competency questions, and expert review (3 Kikuyu consultants).

The Neo4j implementation (v5.x) translated OWL to property graphs: 947 nodes, 1,247 edges, average degree 2.6. Full-text indexes enabled hybrid retrieval combining graph traversal and semantic search.

\section{Phase 4: System Development}
\label{sec:phase_modeling}

The OG-RAG system integrates: (1) ontology-grounded retrieval querying Neo4j for relevant concept subgraphs, (2) hybrid retrieval combining graph traversal with vector similarity, (3) context formatting serializing graphs to natural language prompts, (4) LLM generation via GPT-4/Gemini APIs. Baselines: Raw LLM (zero-shot), Traditional RAG (top-3 vector retrieval). Technical details in Chapter~\ref{chap:design}.

All three systems use identical LLM backends (GPT-4-turbo) and generation parameters (temperature=0.3, max\_tokens=150) to isolate the effect of retrieval mechanisms.

\subsection{Prompt Engineering Strategy}
\label{subsec:prompt_engineering}

Prompt templates for the OG-RAG system were iteratively refined through pilot testing on 20 held-out proverbs. The final prompt structure consists of five sections:

\begin{enumerate}
    \item \textbf{Role Definition}: Establishes the LLM's identity as a cultural translation expert specializing in Kikuyu proverbs.
    
    \item \textbf{Task Specification}: Clearly defines the translation objective, emphasizing cultural fidelity over literal accuracy.
    
    \item \textbf{Structured Context}: Presents retrieved ontological knowledge organized by category (Cultural Themes, Usage Contexts, Related Proverbs, Moral Lessons).
    
    \item \textbf{Input Proverb}: The Kikuyu text to translate, along with its literal word-by-word gloss.
    
    \item \textbf{Output Format}: Specifies the desired translation structure (English translation + brief cultural explanation).
\end{enumerate}

Example prompt structure (simplified):

\begin{verbatim}
You are an expert in Kikuyu culture specializing in proverb translation.

CULTURAL CONTEXT (from ontology):
- Theme: Reciprocity, Community
- Usage: Wealth distribution ceremonies, advice to young people
- Related concepts: Ngwatio system, collective prosperity
- Moral lesson: People are the ultimate wealth

PROVERB TO TRANSLATE:
Kikuyu: "Andu ni indo"
Literal: "People are wealth"

Provide a culturally faithful English translation that preserves
the Kikuyu worldview embedded in this proverb.
\end{verbatim}

\section{Phase 5: Evaluation Methodology}
\label{sec:phase_evaluation}

The evaluation phase employed a multi-dimensional framework combining quantitative metrics with qualitative analysis to rigorously assess system performance.

\subsection{Evaluation Dataset}
\label{subsec:eval_dataset}

The evaluation utilized the complete 100-proverb corpus from Ireri's collection, with each proverb translated by all three systems (Raw LLM, Traditional RAG, OG-RAG), yielding 300 total translations for assessment.

To control for ordering effects, proverbs were presented to evaluators in randomized order with system outputs anonymized (labeled System A, B, C). Evaluators were blind to which system generated each translation.

\subsection{Evaluation Metrics}
\label{subsec:eval_metrics}

A two-dimensional metric framework was developed based on preliminary analysis showing that existing MT metrics (BLEU, CHRF++) correlated poorly with expert judgments of cultural fidelity (Pearson's $r = 0.32$, $p < 0.01$).

\subsubsection{Cultural Authenticity Metric}
\label{subsubsec:cultural_authenticity}

Cultural authenticity measures whether the translation accurately conveys the cultural worldview, values, and contextual appropriateness of the original proverb. Evaluators rated translations on a 0-1 continuous scale considering:

\begin{itemize}
    \item \textbf{Thematic Accuracy (40\%)}: Does the translation preserve the core cultural themes (e.g., reciprocity, community, patience)?
    
    \item \textbf{Metaphorical Preservation (30\%)}: Are culturally specific metaphors either preserved or appropriately adapted for English-speaking audiences?
    
    \item \textbf{Contextual Appropriateness (20\%)}: Would the translation be suitable for the cultural contexts where the original is used?
    
    \item \textbf{Value Alignment (10\%)}: Does the translation reflect Kikuyu values without imposing Western cultural assumptions?
\end{itemize}

The weighted composite score provides a single cultural authenticity value per translation.

\subsubsection{Translation Fidelity Metric}
\label{subsubsec:translation_fidelity}

Translation fidelity measures linguistic accuracy and semantic correspondence to the original meaning, independent of cultural depth. Components include:

\begin{itemize}
    \item \textbf{Semantic Accuracy (50\%)}: Does the translation capture the denotative meaning of the Kikuyu text?
    
    \item \textbf{Fluency (30\%)}: Is the English natural and grammatically correct?
    
    \item \textbf{Completeness (20\%)}: Are all elements of the original proverb reflected in the translation?
\end{itemize}

\subsubsection{Overall Quality Metric}
\label{subsubsec:overall_quality}

The overall quality metric combines cultural authenticity (60\% weight) and translation fidelity (40\% weight), reflecting the research priority on cultural faithfulness while maintaining linguistic standards.

\subsection{Evaluation Procedures}
\label{subsec:eval_procedures}

\subsubsection{Expert Human Evaluation}
\label{subsubsec:expert_evaluation}

Two native Kikuyu speakers with demonstrated cultural competence (both elders over age 60 with traditional knowledge) independently evaluated all 300 translations. Each evaluator:

\begin{enumerate}
    \item Read the original Kikuyu proverb
    \item Reviewed the three anonymized translations (randomized order)
    \item Assigned cultural authenticity and translation fidelity scores
    \item Provided qualitative comments on notable strengths/weaknesses
\end{enumerate}

Inter-rater reliability was assessed using Krippendorff's alpha, achieving $\alpha = 0.78$ for cultural authenticity and $\alpha = 0.82$ for translation fidelity, indicating substantial agreement. Discrepancies exceeding 0.3 points were adjudicated through discussion.

\subsubsection{LLM-Assisted Evaluation}
\label{subsubsec:llm_evaluation}

As a supplementary evaluation method, Gemini 2.5 Pro was configured as an automated judge to provide additional assessment perspectives. Cohere's Command-R-Plus was configured as a fallback provider to ensure evaluation continuity in case of API limitations, though the primary evaluation successfully completed using Gemini. The LLM judge evaluated translations along four dimensions:

\begin{itemize}
    \item Cultural Authenticity (0-1 scale)
    \item Translation Accuracy (0-1 scale)
    \item Contextual Appropriateness (0-1 scale)
    \item Overall Quality (0-1 scale)
\end{itemize}

While LLM judgments showed moderate correlation with expert ratings ($r = 0.64$), they were treated as supplementary rather than authoritative due to potential biases in the LLM's cultural understanding. The primary analysis relies on human expert evaluation.

\subsection{Statistical Analysis Methods}
\label{subsec:statistical_methods}

Quantitative comparisons between systems employed:

\paragraph{Paired t-tests} To assess whether mean score differences between systems were statistically significant ($\alpha = 0.05$). Paired tests were appropriate because each proverb was evaluated across all three systems, creating natural pairs.

\paragraph{Effect Size Calculation} Cohen's $d$ was computed to quantify practical significance beyond statistical significance. Effect sizes were interpreted as: small ($d = 0.2$), medium ($d = 0.5$), large ($d = 0.8$).

\paragraph{Distribution Analysis} Kernel density estimation and box plots visualized score distributions to identify systematic patterns (e.g., reduced variance, fewer catastrophic failures).

All statistical analyses were conducted using SciPy (version 1.11) and visualizations created with Matplotlib/Seaborn.

\subsection{Qualitative Analysis Methods}
\label{subsec:qualitative_methods}

Beyond quantitative metrics, qualitative analysis identified patterns and mechanisms:

\paragraph{Exemplar Analysis} High-scoring and low-scoring translations from each system were examined in detail to understand success and failure modes.

\paragraph{Error Categorization} Translation errors were taxonomized into categories (e.g., semantic collapse, cultural drift, metaphor loss) to identify systematic weaknesses.

\paragraph{Mechanism Elicitation} By comparing translations with the retrieved ontological contexts, we identified how specific ontology features (e.g., cultural theme annotations, usage context descriptions) influenced generation quality.

This qualitative analysis complements quantitative metrics by explaining \emph{why} performance differences emerged, not merely \emph{that} they exist.

\subsection{Annotator Information and Inter-Rater Reliability}
\label{subsec:annotator_info}

Two expert evaluators were recruited through the African Proverbs Working Group: both are L1 Kikuyu speakers (ages 62 and 68 from Nyeri and Murang'a counties), recognized community elders with expertise in traditional oral literature, and experienced in translating Kikuyu cultural materials.

\begin{table}[h]
\centering
\begin{tabular}{lcc}
\hline
\textbf{Characteristic} & \textbf{Evaluator 1} & \textbf{Evaluator 2} \\
\hline
Age & 62 years & 68 years \\
Primary Dialect & Nyeri & Murang'a \\
Education & Secondary & University (BA) \\
Translation Projects & 5+ & 10+ \\
\hline
\end{tabular}
\caption{Evaluator demographics}
\label{tab:evaluator_demographics}
\end{table}

Inter-rater reliability (IRR) was assessed using Krippendorff's alpha, showing substantial agreement: Cultural Authenticity $\alpha = 0.78$, Translation Fidelity $\alpha = 0.82$, Overall Quality $\alpha = 0.80$. All values exceed the $\alpha \geq 0.67$ threshold for tentative conclusions \cite{krippendorff2018content}. For 18 proverbs (6\%) where scores diverged by $>0.3$ points, structured adjudication sessions achieved consensus. No systematic bias was detected (mean difference $<0.02$ points, $p > 0.4$).

\subsubsection{Scoring Rubrics}
\label{subsubsec:scoring_rubrics}

To standardize evaluation and enhance reproducibility, detailed scoring rubrics were provided to evaluators. Tables~\ref{tab:cultural_rubric},~\ref{tab:fidelity_rubric}, and~\ref{tab:eval_scenarios} present these rubrics and example applications.

\begin{table}[p]
\centering
\small
\begin{tabular}{p{1.5cm}p{5cm}p{6cm}}
\hline
\textbf{Score Range} & \textbf{Criteria} & \textbf{Indicators} \\
\hline
0.9--1.0 & \textbf{Exemplary Cultural Authenticity} & Translation fully captures cultural themes, metaphors, and values. Would be recognized by Kikuyu speakers as culturally authentic. Appropriate for all traditional usage contexts. \\
\hline
0.7--0.89 & \textbf{Strong Cultural Grounding} & Core cultural themes preserved with minor simplifications. Metaphors adapted appropriately for English audiences. Some contextual nuance lost but values intact. \\
\hline
0.5--0.69 & \textbf{Moderate Cultural Accuracy} & Main theme recognizable but culturally specific details diluted. Metaphors partially preserved or generalized. Suitable for general audiences but lacks depth for cultural insiders. \\
\hline
0.3--0.49 & \textbf{Limited Cultural Fidelity} & Generic interpretation missing cultural specificity. Metaphors literalized or replaced with Western equivalents. Cultural values implied but not explicit. \\
\hline
0.0--0.29 & \textbf{Culturally Inappropriate} & Translation fundamentally misrepresents cultural meaning. Metaphors nonsensical or offensive. Imposes Western cultural frames. Unusable for cultural education. \\
\hline
\end{tabular}
\caption{Cultural Authenticity Scoring Rubric (0-1 continuous scale)}
\label{tab:cultural_rubric}
\end{table}

\begin{table}[p]
\centering
\small
\begin{tabular}{p{1.5cm}p{5cm}p{6cm}}
\hline
\textbf{Score Range} & \textbf{Criteria} & \textbf{Indicators} \\
\hline
0.9--1.0 & \textbf{Excellent Translation Quality} & Semantically accurate, fluent English, complete coverage of source text. No omissions or additions. Reads naturally to native English speakers. \\
\hline
0.7--0.89 & \textbf{Good Translation Quality} & Meaning preserved with minor infelicities. Fluent but may contain slight awkwardness. One or two minor omissions acceptable if meaning retained. \\
\hline
0.5--0.69 & \textbf{Adequate Translation Quality} & Core meaning comprehensible but some semantic drift. Moderately fluent with noticeable non-native phrasing. Some information loss tolerable. \\
\hline
0.3--0.49 & \textbf{Poor Translation Quality} & Significant semantic errors or omissions. Choppy or unnatural English. Meaning partially obscured. Requires effort to understand. \\
\hline
0.0--0.29 & \textbf{Unacceptable Translation Quality} & Severe semantic errors, unintelligible English, or major omissions. Translation fails basic communication function. \\
\hline
\end{tabular}
\caption{Translation Fidelity Scoring Rubric (0-1 continuous scale)}
\label{tab:fidelity_rubric}
\end{table}

\begin{table}[p]
\centering
\small
\begin{tabular}{p{0.5cm}p{4cm}p{5cm}p{3cm}}
\hline
\textbf{ID} & \textbf{Kikuyu Proverb} & \textbf{Cultural Theme} & \textbf{Evaluation Focus} \\
\hline
P1 & \emph{Andu ni indo} & Community as wealth; reciprocity over material accumulation & Tests cultural value translation vs. literal rendering \\
\hline
P2 & \emph{Mũtumia mũgima ndagagĩrwo nĩ thĩĩna} & Gender roles; holistic prosperity concept & Tests metaphorical preservation and cultural appropriateness \\
\hline
P3 & \emph{Aikaragia mbia ta njuu ngigi} & Vigilance metaphor from stork hunting behavior & Tests handling of culturally specific animal symbolism \\
\hline
P4 & \emph{Gũtirĩ mũndũ ũthiiaga akorwo na mũno} & Continuous improvement; anti-complacency & Tests abstract value communication \\
\hline
P5 & \emph{Ũgima ũrĩa ũtarĩ ũrũgamĩro nĩ ta mũtĩ ũtarĩ mĩri} & Rootedness metaphor for sustainable wealth & Tests complex metaphorical mapping \\
\hline
\end{tabular}
\caption{Representative Evaluation Scenarios}
\label{tab:eval_scenarios}
\end{table}

These rubrics operationalize the abstract concepts of "cultural authenticity" and "translation fidelity" into concrete, gradable criteria. During evaluator training, we discussed example translations at each rubric level to calibrate shared understanding before independent evaluation commenced.

\section{Phase 6: Deployment Considerations}
\label{sec:phase_deployment}

While full production deployment is beyond this thesis scope, we address key scalability and ethical considerations for future application. Scaling the 100-proverb ontology to comprehensive coverage (3,000+ proverbs) would require semi-automated concept extraction and collaborative editing interfaces. Query optimization (current latency: 2.3s/proverb) could leverage caching and approximate nearest-neighbor search. Ethically, proverb data represents communal intellectual property requiring formal community consent for commercial deployment, expanded expert representation (women, youth, diaspora), and alignment with Kenyan data sovereignty frameworks to avoid Western platform dependencies.

\begin{enumerate}
    \item \textbf{Multilingual Extension}: Adapting the OG-RAG approach to other Bantu languages (Kamba, Gusii, Luhya) to test generalizability.
    
    \item \textbf{Bidirectional Translation}: Developing English-to-Kikuyu capabilities, which require different cultural adaptation strategies.
    
    \item \textbf{Ontology Evolution}: Mechanisms for capturing contemporary proverb reinterpretations and newly coined sayings.
    
    \item \textbf{Educational Applications}: Integrating the system into Kikuyu language learning platforms as a cultural education tool.
    
    \item \textbf{Participatory Ontology Development}: Creating community-driven platforms where Kikuyu speakers worldwide can contribute proverbs and cultural annotations.
\end{enumerate}

These directions position the current work as an initial prototype within a longer-term research program on ontology-grounded cultural AI.

\section{Methodological Limitations}
\label{sec:methodological_limitations}

We prioritized depth over breadth (100 wealth-domain proverbs vs. 3,000+ corpus-wide), cultural validity over scale (expert judgment vs. automated metrics), and ontological coherence over completeness. Limitations include: (1) \textbf{Sample constraints}: findings specific to wealth domain, requiring validation for kinship, agriculture, governance themes; (2) \textbf{Ontology gaps}: 12\% of proverbs invoked concepts not formally represented; (3) \textbf{Evaluation subjectivity}: IRR $\alpha > 0.75$ substantial but evaluator variance unavoidable; (4) \textbf{Temporal validity}: captures traditional interpretations, may miss contemporary urban reinterpretations; (5) \textbf{Generalization boundaries}: methodology transferable to other Bantu languages but results not automatically generalizable. These boundaries establish confidence limits while suggesting concrete expansion directions.

\section{Summary}
\label{sec:methodology_summary}

This chapter presented a CRISP-DM-based methodology for ontology-driven cultural NLP research. The six-phase approach (problem definition, data understanding, ontology construction with 847 concepts, system development, multi-dimensional evaluation, deployment considerations) demonstrates systematic elicitation and operationalization of tacit cultural knowledge. Key contributions: (1) ontology construction process for cultural knowledge representation, (2) two-dimensional evaluation framework (cultural authenticity + translation fidelity) addressing gaps in culturally sensitive translation assessment. Acknowledged limitations establish confidence bounds while suggesting expansion directions. Chapters~\ref{chap:design} and~\ref{ch:evaluation} detail technical implementation and empirical results.
